\scanpage{105}
\mspage{3.4}
% archive sheet number: 104
unramified over $S$. Since $\overline D=\overline R$, we have a natural epimorphism
$D\mathbin{\ooalign{\hfil$\mathrm{x}$\hfil\cr\hfil$\handwritten{\bigcirc}$\hfil\cr}}_R S\xrightarrow{\handwritten{\varphi}}\overline S$. Let $\underline M=\mathrm{Ker}\,\handwritten{\varphi}$ and set $\overset{\handwritten{\sim}}{S}=(D\mathbin{\ooalign{\hfil$\mathrm{x}$\hfil\cr\hfil$\handwritten{\bigcirc}$\hfil\cr}}_R S)_{\underline M}$. Then $\overset{\handwritten{\sim}}{S}$ is flat and
unramified over $S$. Therefore $e(S/R)=e(\overset{\handwritten{\sim}}{S}/D)$ by lemma 2.3. From
lemma 3.6, we get $e(\overset{\handwritten{\sim}}{S}/D)=\handwritten{\sum}_{\underline p}L(\overset{\handwritten{\sim}}{S}_{\underline p})e(\overset{\handwritten{\sim}}{S}/\underline p\mid D)$ where $\underline p$ ranges through the
minimal primes in $\overset{\handwritten{\sim}}{S}$ such that $\dim \overset{\handwritten{\sim}}{S}/\underline p=\dim \overset{\handwritten{\sim}}{S}$. Now if we let $\underline p_0={}$
$\mathrm{Ker}(D\mathbin{\ooalign{\hfil$\mathrm{x}$\hfil\cr\hfil$\handwritten{\bigcirc}$\hfil\cr}}_R S\longrightarrow DS)$, then $\underline p_0\cap R=0$ and thus we are in a position to
apply lemma 4.5. Thus we have that $e(\overset{\handwritten{\sim}}{S}/D)=\handwritten{\sum}_{\underline p}L(\overset{\handwritten{\sim}}{S}_{\underline p})e(\overset{\handwritten{\sim}}{S}/\underline p\mid D)$ where
$\underline p\cap R=0$. We claim that the local $D$-algebras $\overset{\handwritten{\sim}}{S}/\underline p$ are all finite $D$-module.
For, since there is one-one correspondence between the prime ideals $\underline p$ in
$D\mathbin{\ooalign{\hfil$\mathrm{x}$\hfil\cr\hfil$\handwritten{\bigcirc}$\hfil\cr}}_R S$ with $\underline p\cap R=0$ and the prime ideals in $(D\mathbin{\ooalign{\hfil$\mathrm{x}$\hfil\cr\hfil$\handwritten{\bigcirc}$\hfil\cr}}_R S)\mathbin{\ooalign{\hfil$\mathrm{x}$\hfil\cr\hfil$\handwritten{\bigcirc}$\hfil\cr}}_R k=L'_D\mathbin{\ooalign{\hfil$\mathrm{x}$\hfil\cr\hfil$\handwritten{\bigcirc}$\hfil\cr}}_k K$
(where $L'_D$ is the quotient field of $D$), we see that $\overset{\handwritten{\sim}}{S}/\underline p$ is isomorphic to (and
hence we shall identify) a localization of $\handwritten{\sigma}(S)D$ as $D$-algebras, for some
conjugation $\handwritten{\sigma}$ of $S$ over\unsure{'} $R$. (The reader may observe that $\handwritten{\sigma}(S)D$ corresponding
to $\overset{\handwritten{\sim}}{S}/\underline p$ has, by the definition of $\overset{\handwritten{\sim}}{S}=(D\mathbin{\ooalign{\hfil$\mathrm{x}$\hfil\cr\hfil$\handwritten{\bigcirc}$\hfil\cr}}_R S)_{\underline M}$, a maximal ideal which
dominates both the maximal ideals of $S^{\handwritten{\sigma}}$ and $D$). The quotient field of
$\overset{\handwritten{\sim}}{S}/\underline p$ is $L'_D\handwritten{\sigma}(K)$ and $L'_D\handwritten{\sigma}(K')$ is the separable closure of $L'_D$ in
$L'_D\handwritten{\sigma}(K)$. Since $\overset{\handwritten{\sim}}{S}/\underline p$ is a local $D$-algebra, we can find a local $D$-algebra
$A$ with quotient field $L'_D\handwritten{\sigma}(K')$ such that $D\handwritten{\subset} A\handwritten{\subset} \overset{\handwritten{\sim}}{S}/\underline p$. Since $L'_D\handwritten{\sigma}(K)$
is purely inseparable over $L'_D\handwritten{\sigma}(K')$, every integral extension of a local
ring $A$ contained in $L'_D\handwritten{\sigma}(K)$ is local and consequently local $A$-algebra
$\overset{\handwritten{\sim}}{S}/\underline p$ is a finite $A$-module. On the other hand, since the integral closure
of $D$ in $L'$ is a finite local $D$-algebra (proposition 1.1), every local $D$-algebra
contained in $L'$ is a finite local $D$-algebra. In particular, the
local $D$-algebra $A$, being contained in $L'_D\handwritten{\sigma}(K')\handwritten{\subset} L$, is a finite local
$D$-algebra. Therefore we conclude that $\overset{\handwritten{\sim}}{S}/\underline p$ is a finite local $D$-algebra
for all $\underline p$ with $\underline p\cap R=0$. Furthermore the residue field of $S/\underline p$, being
isomorphic to $\overline S$, is a simple extension of $\overline D=\overline R$ by the hypothesis. It
